\section{符号说明}

本文中使用的主要符号及其含义如表~\ref{tab:symbols}~所示。

\begingroup
\small
\setlength{\tabcolsep}{3pt}
\begin{longtable}{p{1.6cm}p{4.0cm}p{6.5cm}}
\caption{主要符号说明}\label{tab:symbols}\\
\toprule
符号 & 含义 & 单位/说明 \\
\midrule
\endfirsthead
\toprule
符号 & 含义 & 单位/说明 \\
\midrule
\endhead
$t$ & 时间（年份） & 数据年份 \\
$T$ & 年份向量 & $T=(1992,2002,2012,2020,2023)$ \\
$Y$ & 超重肥胖率观测值向量 & $Y=(20.0,29.9,42.0,50.7,57.0)$，单位\% \\
$\Delta t_k$ & 相邻观测的时间间隔 & $\Delta t=(10,10,8,3)$年 \\
$x^{(1)}$ & 一次累加生成序列（AGO） & \\
$a$ & GM(1,1)模型的发展系数 & \\
$u$ & GM(1,1)模型的灰作用量 & \\
$K$ & Logistic/Gompertz模型的环境容纳量 & \\
$r$ & Logistic模型的内在增长率 & \\
$t_0$ & Logistic模型的曲线拐点（年） & \\
$b$ & Gompertz模型的位置参数 & \\
$\mathrm{RMSE}$ & 均方根误差 & \\
$\mathrm{MAPE}$ & 平均绝对百分比误差 & \\
$R^2$ & 拟合优度（决定系数） & \\
\midrule
$x_{ij}$ & 第$i$项减重维持因素的第$j$个原始指标值 & 问题二决策矩阵元素 \\
$z_{ij}$ & 原始指标的标准化值 & 分别用于CRITIC极差标准化和TOPSIS向量标准化 \\
$\sigma_j$ & 第$j$个指标标准化值的标准差 & 衡量指标的区分能力 \\
$r_{jk}$ & 第$j$、$k$个指标之间的相关系数 & 衡量指标信息的重复程度 \\
$C_j$ & 第$j$个指标的CRITIC信息量 & 综合离散程度与冲突程度 \\
$w_j$ & 第$j$个指标的客观权重 & $\sum_j w_j=1$ \\
$v_{ij}$ & 第$i$项因素在第$j$个指标上的加权标准化值 & TOPSIS加权矩阵元素 \\
$A^+$ & 正理想解 & 各效益型指标的最优值集合 \\
$A^-$ & 负理想解 & 各效益型指标的最差值集合 \\
$D_i^+$ & 第$i$项因素到正理想解的距离 & 越小越接近最优证据组合 \\
$D_i^-$ & 第$i$项因素到负理想解的距离 & 越大越远离最差证据组合 \\
$S_i$ & 第$i$项因素的TOPSIS贴近度 & $0\leq S_i\leq1$，越大优先级越高 \\
\bottomrule
\end{longtable}
\endgroup

问题一使用的5组历史数据时间间隔并不均匀（$\Delta t = (10, 10, 8, 3)$年），因此在模型建立时需特别考虑非等间距的影响；问题二的三个评价指标均统一为效益型指标。

% ==================== 五、模型的建立与求解 ====================
